\subsection{问题、记号与主结论}

全文在复数域 \(\C\) 上工作。令
\[
  V=\Span\{x_{ij}:0\leq i,j\leq3\},\qquad \dim V=16.
\]

\begin{definition}
四次齐次多项式 \(f\in\Sym^4V\) 的 \emph{Chow 秩}是最小整数 \(r\)，
使得存在 \(\ell_{st}\in V\) 满足
\[
  f=\sum_{s=1}^{r}\ell_{s1}\ell_{s2}\ell_{s3}\ell_{s4}.
\]
每个乘积 \(\ell_1\ell_2\ell_3\ell_4\) 称为一个 Chow 项。
\end{definition}

\begin{theorem}[主定理]\label{n4:thm:main}
\[
  \boxed{\ChowRank(\perm_4)=8}.
\]
\end{theorem}

证明分为两部分。第~\ref{n4:sec:upper} 节给出八项分解；其余各节建立
不存在七项分解的下界。

\subsection{二阶导数空间与 Koszul 展平}\label{n4:sec:koszul}

对 \(f\in\Sym^4V\)，记其中间催化映射为
\[
  C_{2,2}(f):\Sym^2V^*\longrightarrow\Sym^2V,
  \qquad D\longmapsto D(f),
\]
并记
\[
  E_f:=\im C_{2,2}(f)\subseteq\Sym^2V.
\]
定义 Koszul 微分
\[
  \delta:\Sym^2V\otimes V\longrightarrow V\otimes\bigwedge\nolimits^2V
\]
为
\begin{equation}\label{n4:eq:delta}
  \delta\bigl((uv)\otimes w\bigr)
  =v\otimes(u\wedge w)+u\otimes(v\wedge w),
\end{equation}
并线性延拓。本文使用的展平是
\[
  K(f):=\delta\circ\bigl(C_{2,2}(f)\otimes\id_V\bigr).
\]
因此
\begin{equation}\label{n4:eq:image-k}
  \im K(f)=\delta(E_f\otimes V).
\end{equation}

在单项式基下，中间催化矩阵是对称的。对上式取转置可得
\begin{equation}\label{n4:eq:row-containment}
  \row K(f)\subseteq E_f\otimes V^*.
\end{equation}
这一行空间包含关系将在双商估计中使用。

\subsubsection{四阶永久多项式的二阶导数空间}

记 \(E:=E_{\perm_4}\)。直接求二阶导数得到
\begin{equation}\label{n4:eq:E-basis}
  E=\Span\left\{
    x_{ia}x_{jb}+x_{ib}x_{ja}:
    0\leq i<j\leq3,\ 0\leq a<b\leq3
  \right\}.
\end{equation}
这些二次型具有互不混淆的行对、列对支撑，故线性无关。因此
\[
  \dim E=\binom42^2=36.
\]

\begin{proposition}\label{n4:prop:base-ranks}
有
\[
  \rank K(\perm_4)=560.
\]
若 \(T=\ell_1\ell_2\ell_3\ell_4\) 是任意 Chow 项，则
\[
  \rank K(T)\leq92;
\]
当 \(\ell_1,\ldots,\ell_4\) 线性无关时等号成立。
\end{proposition}

\begin{proof}
先看 \(\perm_4\)。映射
\[
  \left.\delta\right|_{E\otimes V}:E\otimes V\longrightarrow
  V\otimes\bigwedge\nolimits^2V
\]
的定义域维数为 \(36\cdot16=576\)。若三次型 \(g\in\Sym^3V\)
的所有一阶偏导都属于 \(E\)，定义
\[
  \iota(g):=\sum_{a=0}^{15}
  \frac{\partial g}{\partial x_a}\otimes x_a\in E\otimes V.
\]
把式~\eqref{n4:eq:delta} 逐个用于三次单项式，可见所有楔积项成对
相消，故 \(\delta(\iota(g))=0\)。另一方面，张量乘法映射将
\(\iota(g)\) 送到 Euler 恒等式
\[
  \sum_{a=0}^{15}x_a\frac{\partial g}{\partial x_a}=3g.
\]
因此 \(g\mapsto\iota(g)\) 是单射。

对每个三行子集
\(I\subset\{0,1,2,3\}\) 和三列子集 \(J\subset\{0,1,2,3\}\)，令
\[
  p_{I,J}:=\perm(x_{ij})_{i\in I,j\in J}.
\]
共有 \(\binom43^2=16\) 个这样的三阶子永久式。它们支撑不同，故线性
无关；而每个一阶导数 \(\partial p_{I,J}/\partial x_{ij}\) 都属于
\(E\)。由上面的单射，它们给出
\(\delta|_{E\otimes V}\) 的十六个线性无关核向量。于是
\[
  \dim\ker(\delta|_{E\otimes V})\geq16,\qquad
  \rank K(\perm_4)\leq576-16=560.
\]
反方向，在式~\eqref{n4:eq:E-basis}、\(V\) 的坐标基和
\(V\otimes\bigwedge^2V\) 的字典序基下，对整数矩阵
\(\delta|_{E\otimes V}\) 作消元，得到一个在
\(\F_{1000003}\) 上非零的 \(560\times560\) 子式。因此同一个整数
子式的行列式不是零整数，从而在特征零中
\(\rank K(\perm_4)\geq560\)。两边合并即得等号。

再看线性无关的 Chow 项。经 \(GL(V)\) 变换可令
\[
  T=x_0x_1x_2x_3.
\]
此时
\[
  E_T=\Span\{x_ix_j:0\leq i<j\leq3\},\qquad \dim E_T=6,
\]
所以定义域维数为 \(6\cdot16=96\)。四个三次单项式
\[
  x_0x_1x_2,\quad x_0x_1x_3,\quad
  x_0x_2x_3,\quad x_1x_2x_3
\]
的所有一阶偏导都属于 \(E_T\)，故同一个映射 \(\iota\) 给出四个
线性无关的核向量。因此 \(\rank K(T)\leq96-4=92\)。
同一整数矩阵在 \(\F_{1000003}\) 上存在非零的 \(92\times92\)
子式，故 \(\rank K(T)=92\)。

最后，线性无关的四元组在参数空间 \((V^*)^4\) 中构成稠密开集，
且所有这些四元组经 \(GL(V)\) 变换后都有秩 \(92\)。矩阵秩大于
\(92\) 的点集是开集；若某个退化四元组的秩大于 \(92\)，该开集
会与上述稠密开集相交，产生矛盾。因此任意 Chow 项的秩均不超过
\(92\)。
\end{proof}

\begin{remark}
上述模素数计算是严格证书而非概率证据：一个整数子式模某个素数
非零，立即推出该子式在 \(\mathbb Z\) 及 \(\C\) 上非零。第
~\ref{n4:sec:certificate} 节给出独立重建程序和输出。
\end{remark}

\subsection{低秩二次型的分类}\label{n4:sec:classification}

本节完全是手工线性代数，不依赖计算机搜索。

用六个行对和六个列对给式~\eqref{n4:eq:E-basis} 中的基编号。任意
\(q\in E\) 对应一个 \(6\times6\) 系数矩阵 \(C\)。对固定行对
\(i<j\)，把 \(C\) 的相应一行写成对称零对角 \(4\times4\) 矩阵
\[
  D_{ij}[a,b]=C[(i,j),(a,b)]\quad(a\neq b),\qquad
  D_{ij}[a,a]=0.
\]
将 \(q\) 的变量按矩阵的四行分组，则 \(q\) 的对称系数矩阵为
\begin{equation}\label{n4:eq:block-matrix}
  M(q)=(D_{ij})_{0\leq i,j\leq3},\qquad
  D_{ji}=D_{ij},\quad D_{ii}=0.
\end{equation}
采用 Hessian 还是对称系数矩阵只相差非零标量，不影响秩。

\begin{lemma}\label{n4:lem:zero-diagonal}
设 \(D\neq0\) 是复数域上的对称零对角矩阵。
\begin{enumerate}[label=(\roman*)]
  \item \(\rank D\neq1\)；
  \item 若 \(\rank D=2\)，且 \(H\) 也是对称零对角矩阵并满足
  \(\im H\subseteq\im D\)，则 \(H\) 是 \(D\) 的标量倍数。
\end{enumerate}
\end{lemma}

\begin{proof}
若对称矩阵的秩为 \(1\)，则它可写成 \(\lambda uu^{\mathsf T}\)。
零对角给出 \(\lambda u_i^2=0\) 对每个 \(i\) 成立，故 \(u=0\)，
与 \(D\neq0\) 矛盾。

若 \(\rank D=2\)，则在 \(\C\) 上可取线性无关向量 \(a,b\)，使
\[
  D=ab^{\mathsf T}+ba^{\mathsf T}.
\]
由零对角得 \(a_kb_k=0\) 对每个 \(k\) 成立，所以 \(a,b\) 的
支撑不交。由于 \(\im H\subseteq\Span(a,b)\) 且 \(H\) 对称，
存在对称 \(2\times2\) 矩阵 \(R\) 使
\[
  H=[\,a\ b\,]R[\,a\ b\,]^{\mathsf T}.
\]
在 \(a\) 与 \(b\) 的非空支撑上分别考察对角元，得到
\(R_{11}=R_{22}=0\)。因此 \(R\) 只剩一个非对角参数，
即 \(H\) 是 \(D\) 的标量倍数。
\end{proof}

\begin{theorem}\label{n4:thm:rank-four-classification}
若 \(0\neq q\in E\) 且 \(\rank M(q)\leq4\)，则
\[
  M(q)=U\otimes W,
\]
其中 \(U,W\) 都是秩为 \(2\) 的对称零对角 \(4\times4\) 矩阵。
特别地，\(\rank M(q)=4\)。
\end{theorem}

\begin{proof}
选取一个非零块 \(D:=D_{ij}\)。对应的主子矩阵为
\[
  N=\begin{bmatrix}0&D\\D&0\end{bmatrix},
  \qquad \rank N=2\rank D\leq\rank M(q)\leq4.
\]
由引理~\ref{n4:lem:zero-diagonal}，非零 \(D\) 不可能秩为 \(1\)，
故 \(\rank D=2\)，并且
\[
  \rank N=4=\rank M(q).
\]
因此 \(N\) 已经承担 \(M(q)\) 的全部秩；其余块列都在 \(N\) 的
列空间中。特别地，对每个 \(k\)，
\[
  \im D_{ik},\ \im D_{jk}\subseteq\im D.
\]
由引理~\ref{n4:lem:zero-diagonal}(ii)，存在标量 \(a_k,b_k\) 使
\[
  D_{ik}=a_kD,\qquad D_{jk}=b_kD.
\]

取 \(D\) 的任意广义逆 \(G\)，满足 \(DGD=D\)，并令
\[
  N^-=\begin{bmatrix}0&G\\G&0\end{bmatrix}.
\]
由于总矩阵与主块 \(N\) 秩相同，广义 Schur 补必须为零。于是对
任意剩余指标 \(k,l\)，
\begin{align*}
  D_{kl}
  &=
  \begin{bmatrix}D_{ik}\\D_{jk}\end{bmatrix}^{\mathsf T}
  N^-
  \begin{bmatrix}D_{il}\\D_{jl}\end{bmatrix}  \\
  &=(a_kb_l+b_ka_l)D.
\end{align*}
所以所有 \(D_{kl}\) 都是同一个 \(D\) 的标量倍数，因而六乘六
系数矩阵 \(C\) 的秩为 \(1\)。把这些标量写成一个对称零对角矩阵
\(U\)，并记 \(W=D\)，便有
\[
  M(q)=U\otimes W.
\]
最后
\[
  4=\rank M(q)=\rank U\cdot\rank W.
\]
两个非零对称零对角因子都不可能秩为 \(1\)，故二者秩都为 \(2\)。
\end{proof}

\begin{corollary}\label{n4:cor:intersection}
对任意 \(\dim L\leq4\) 的线性子空间 \(L\subseteq V\)，有
\[
  \dim\bigl(E\cap\Sym^2L\bigr)\leq1.
\]
因而，对任意 Chow 项 \(T=\ell_1\ell_2\ell_3\ell_4\)，若
\(F:=E_T\)，则
\[
  \dim(E\cap F)\leq1.
\]
\end{corollary}

\begin{proof}
若 \(0\neq q\in E\cap\Sym^2L\)，则 \(q\) 的矩阵秩至多为
\(\dim L\leq4\)。由定理~\ref{n4:thm:rank-four-classification}，
它的秩恰为 \(4\)，且其本质变量空间为
\[
  \im U\otimes\im W,
\]
维数为 \(4\)。因此必有 \(\dim L=4\)，且该本质变量空间等于
\(L\)。

张量积子空间的因子是唯一的：若非零二维子空间满足
\(S_1\otimes T_1=S_2\otimes T_2\)，则
\(S_1=S_2,\ T_1=T_2\)。而在固定二维像空间中，引理
~\ref{n4:lem:zero-diagonal}(ii) 又说明对称零对角秩二矩阵在相差
标量的意义下唯一。因此，\(E\cap\Sym^2L\) 中任意两个非零元素
必成比例，证明第一式。

对 Chow 项，所有二阶导数都只使用
\(L=\Span(\ell_1,\ldots,\ell_4)\) 中的变量，即
\[
  F\subseteq\Sym^2L.
\]
第二式立即由第一式得到，包括因子重复或线性相关的退化情形。
\end{proof}

\subsection{精确图表证书}\label{n4:sec:chart}

这是证明中唯一实质性的计算机辅助步骤。先定义对每个
\(0\neq v\in V\) 的线性映射
\begin{equation}\label{n4:eq:psi}
  \psi_v:\Sym^2V/E\longrightarrow
  \coker\bigl(\delta(E\otimes V)\bigr),\qquad
  [q]\longmapsto[\delta(q\otimes v)].
\end{equation}
源空间维数为
\[
  \dim\Sym^2V-\dim E=\binom{17}{2}-36=136-36=100.
\]
又由式~\eqref{n4:eq:delta}，
\[
  \delta(v^2\otimes v)=0.
\]
定理~\ref{n4:thm:rank-four-classification} 表明 \(E\) 中非零二次型
的秩至少为 \(4\)，而 \(v^2\) 的秩为 \(1\)，故 \(v^2\notin E\)。
因此
\begin{equation}\label{n4:eq:obvious-kernel}
  0\neq[v^2]\in\ker\psi_v,\qquad \rank\psi_v\leq99.
\end{equation}

\begin{proposition}[图表证书]\label{n4:prop:chart}
对每个 \(0\neq v\in V\)，
\[
  \rank\psi_v=99,\qquad
  \ker\psi_v=\Span([v^2]).
\]
\end{proposition}

\begin{proof}
先考虑仿射图表 \(v_{00}\neq0\)。固定如下字典序基：
\[
  \Sym^2V:\ x_ax_b\ (0\leq a\leq b<16),\qquad
  V\otimes\bigwedge\nolimits^2V:\
  x_a\otimes(x_b\wedge x_c)\ (b<c).
\]
程序先从 \(\delta(E\otimes V)\) 中选取 \(560\) 个独立列及其
\(560\) 个主元行，再从 \(\Sym^2V/E\) 中确定 \(99\) 个二次类，
并选取额外 \(99\) 行。对
\[
  v=\sum_{k=0}^{15}v_kx_k,\qquad v_0=v_{00},
\]
所得 \(659\times659\) 子矩阵具有分块形式
\begin{equation}\label{n4:eq:659}
  \mathcal M(v)=
  \begin{bmatrix}
    A_0&Q_0(v)\\
    A_1&Q_1(v)
  \end{bmatrix},
\end{equation}
其中 \(A_0\) 为 \(560\times560\)，而右侧有 \(99\) 列。
所有条目均为整数线性式。精确有理数消元得到
\[
  \det A_0=1.
\]
令
\[
  S(v):=Q_1(v)-A_1A_0^{-1}Q_0(v)
       =\sum_{k=0}^{15}v_kS_k.
\]
直接在 \(\Q\) 上验算得到
\begin{equation}\label{n4:eq:S0}
  \det S_0=-32768.
\end{equation}
对 \(k=1,\ldots,15\) 定义 \(N_k=S_0^{-1}S_k\)。证书逐项验证：
\begin{enumerate}[label=(\arabic*)]
  \item 所有 \(N_k\) 的对角元均为零；
  \item 这十五个矩阵的非零位置并集共有 \(156\) 个位置；
  \item 以非零位置 \((a,b)\) 为有向边 \(a\to b\) 所得并图无环。
\end{enumerate}
第三点给出一个公共拓扑序；按此顺序同时置换行列后，每个 \(N_k\)
都严格上三角。因此
\begin{align}
  \det\mathcal M(v)
  &=\det A_0\cdot\det S(v)\notag\\
  &=\det A_0\cdot\det S_0\cdot
    \det\left(v_0I_{99}+\sum_{k=1}^{15}v_kN_k\right)\notag\\
  &=-32768\,v_{00}^{99}.\label{n4:eq:chart-determinant}
\end{align}
当 \(v_{00}\neq0\) 时该子式非零，所以 \(\rank\psi_v\geq99\)。
结合式~\eqref{n4:eq:obvious-kernel}，得到
\[
  \rank\psi_v=99,\qquad
  \ker\psi_v=\Span([v^2]).
\]

最后，\(S_4\times S_4\) 分别置换变量矩阵的行和列。该群保持
\(\perm_4\)、\(E\) 与 Koszul 微分，并在十六个坐标上可迁。
任取 \(v\neq0\)，至少一个坐标非零；用群作用将该坐标移到
\((0,0)\)，便归约到已经证明的图表。因此结论对所有 \(v\neq0\)
成立。
\end{proof}

\begin{remark}
式~\eqref{n4:eq:chart-determinant} 不是数值拟合：矩阵
\(A_0,S_0,N_k\) 均在 \(\Q\) 上构造和消元，公共三角化通过有限
有向图的无环性逐边核验。模素数只用于确定候选行列索引；一旦索引
固定，最终行列式恒等式完全在 \(\Q\) 上重新计算。
\end{remark}

\subsection{一个新二次方向至少增加十五维}\label{n4:sec:extension}

\begin{lemma}[扩张引理]\label{n4:lem:extension}
若 \(q\in\Sym^2V\setminus E\)，则
\[
  \rank\delta\bigl((E+\Span(q))\otimes V\bigr)\geq575.
\]
\end{lemma}

\begin{proof}
固定 \(q\notin E\)，定义
\[
  \theta_q:V\longrightarrow
  \coker\bigl(\delta(E\otimes V)\bigr),\qquad
  v\longmapsto\psi_v([q]).
\]
这是关于 \(v\) 的线性映射。

假设存在两个线性无关的 \(v,w\in\ker\theta_q\)。因为
\([q]\neq0\)，命题~\ref{n4:prop:chart} 给出非零标量 \(a,b\) 使
\[
  [q]=a[v^2]=b[w^2]\quad\text{于 }\Sym^2V/E.
\]
从而
\[
  av^2-bw^2\in E.
\]
由于 \(v,w\) 线性无关且 \(a,b\neq0\)，左端是一个非零且秩至多
为 \(2\) 的二次型；这与定理~\ref{n4:thm:rank-four-classification}
推出的 \(E\) 中非零元素秩至少为 \(4\) 矛盾。因此
\[
  \dim\ker\theta_q\leq1,\qquad \rank\theta_q\geq15.
\]
映射 \(\theta_q\) 的像正是加入 \(q\otimes V\) 后在
\(\coker(\delta(E\otimes V))\) 中新增的像。因此由命题
~\ref{n4:prop:base-ranks}，
\[
  \rank\delta\bigl((E+\Span(q))\otimes V\bigr)
  \geq560+15=575.
\]
\end{proof}

\subsection{双商秩不等式}\label{n4:sec:double}

\begin{lemma}[双商秩不等式]\label{n4:lem:double}
对任意两个同型的 \(m\times n\) 矩阵 \(A,B\)，有
\begin{equation}\label{n4:eq:double}
  \rank(A-B)\geq
  \rank[\,A\ B\,]
  +\rank\begin{bmatrix}A\\B\end{bmatrix}
  -\rank A-\rank B.
\end{equation}
\end{lemma}

\begin{proof}
记
\[
  c=\dim(\im A\cap\im B),\qquad
  d=\dim(\row A\cap\row B).
\]
则
\[
  \rank[\,A\ B\,]=\rank A+\rank B-c,
\]
\[
  \rank\begin{bmatrix}A\\B\end{bmatrix}
  =\rank A+\rank B-d.
\]
对 \(x\in\ker(A-B)\)，有 \(Ax=Bx\in\im A\cap\im B\)。线性映射
\[
  \ker(A-B)\longrightarrow\im A\cap\im B,\qquad x\longmapsto Ax
\]
的核恰为 \(\ker A\cap\ker B\)，故
\[
  \dim\ker(A-B)\leq c+\dim(\ker A\cap\ker B).
\]
而
\[
  \dim(\ker A\cap\ker B)
  =n-\rank\begin{bmatrix}A\\B\end{bmatrix}
  =n-\rank A-\rank B+d.
\]
代回并从 \(n\) 中取余维，即得式~\eqref{n4:eq:double}。
\end{proof}

\begin{remark}\label{n4:rem:scaling}
若 \(\alpha\neq0\)，把 \(B\) 换成 \(\alpha B\) 不改变
\([\,A\ B\,]\) 与 \(\bigl[\begin{smallmatrix}A\\B\end{smallmatrix}\bigr]\)
的秩。因此式~\eqref{n4:eq:double} 同样统一控制 \(A-\alpha B\)。
\end{remark}

\subsection{排除七项分解}\label{n4:sec:lower}

\begin{theorem}\label{n4:thm:lower}
\[
  \ChowRank(\perm_4)\geq8.
\]
\end{theorem}

\begin{proof}
固定任意非零 Chow 项 \(T=\ell_1\ell_2\ell_3\ell_4\)，并记
\[
  A=K(\perm_4),\qquad B=K(T),\qquad
  F=E_T,\qquad d=\dim(E\cap F).
\]
由推论~\ref{n4:cor:intersection}，只有 \(d=0\) 或 \(d=1\) 两种
可能。

\medskip
\noindent\textbf{情形一：\(d=0\)。}
由式~\eqref{n4:eq:row-containment}，
\[
  \row A\subseteq E\otimes V^*,\qquad
  \row B\subseteq F\otimes V^*.
\]
且
\[
  (E\otimes V^*)\cap(F\otimes V^*)
  =(E\cap F)\otimes V^*=0.
\]
故 \(\row A\cap\row B=0\)，从而
\[
  \rank\begin{bmatrix}A\\B\end{bmatrix}
  =\rank A+\rank B.
\]
将其代入引理~\ref{n4:lem:double}，并使用
\(\rank[\,A\ B\,]\geq\rank A=560\)，得到对任意
\(\alpha\neq0\)，
\begin{equation}\label{n4:eq:case0}
  \rank(A-\alpha B)\geq560.
\end{equation}

\medskip
\noindent\textbf{情形二：\(d=1\)。}
取 \(0\neq q\in E\cap F\)。因为 \(F\subseteq
\Sym^2\Span(\ell_1,\ldots,\ell_4)\)，且 \(q\) 的秩由定理
~\ref{n4:thm:rank-four-classification} 等于 \(4\)，四个因子必须
张成四维空间，因而线性无关。由命题~\ref{n4:prop:base-ranks}，
\[
  \dim F=6,\qquad \rank B=92.
\]
可取 \(r\in F\setminus E\)。水平拼接矩阵的像包含
\(\delta((E+\Span(r))\otimes V)\)，所以扩张引理给出
\begin{equation}\label{n4:eq:horizontal}
  \rank[\,A\ B\,]\geq575.
\end{equation}
另一方面，由行空间包含关系，
\[
  \dim(\row A\cap\row B)
  \leq\dim\bigl((E\cap F)\otimes V^*\bigr)=16.
\]
因此
\begin{equation}\label{n4:eq:vertical}
  \rank\begin{bmatrix}A\\B\end{bmatrix}
  \geq560+92-16=636.
\end{equation}
把式~\eqref{n4:eq:horizontal} 与式~\eqref{n4:eq:vertical} 代入双商
不等式，并用说明~\ref{n4:rem:scaling}，得到
\begin{equation}\label{n4:eq:case1}
  \rank(A-\alpha B)
  \geq575+636-560-92=559
  \qquad(\alpha\neq0).
\end{equation}

两种情形合并说明：从 \(\perm_4\) 中减去任意一个非零 Chow 项
\(\alpha T\) 后，其 Koszul 展平秩都至少为 \(559\)。

若 \(\perm_4\) 有七项 Chow 分解，选取其中任意一个非零项
\(\alpha T\)，则余项是六个 Chow 项之和。展平关于多项式线性，
矩阵秩又满足次可加性；由命题~\ref{n4:prop:base-ranks}，
\[
  \rank K(\perm_4-\alpha T)
  \leq6\cdot92=552.
\]
这与式~\eqref{n4:eq:case0} 或式~\eqref{n4:eq:case1} 的下界矛盾。
若分解中实际少于七个非零项，则同样有秩至多 \(6\cdot92=552\)，
甚至直接与 \(\rank K(\perm_4)=560\) 矛盾。因此
\(\ChowRank(\perm_4)\geq8\)。
\end{proof}

\subsection{Glynn 恒等式给出的八项上界}\label{n4:sec:upper}

对任意 \(n\)，Glynn 恒等式为
\begin{equation}\label{n4:eq:glynn}
  \perm_n
  =2^{1-n}
  \sum_{\substack{\varepsilon\in\{\pm1\}^n\\\varepsilon_0=1}}
  \left(\prod_{i=0}^{n-1}\varepsilon_i\right)
  \prod_{j=0}^{n-1}
  \left(\sum_{i=0}^{n-1}\varepsilon_i x_{ij}\right).
\end{equation}
求和中共有 \(2^{n-1}\) 个符号向量，每个求和项都是 \(n\) 个
线性形式的乘积。取 \(n=4\)，式~\eqref{n4:eq:glynn} 给出
\(\perm_4\) 的八项 Chow 分解，故
\[
  \ChowRank(\perm_4)\leq8.
\]
与定理~\ref{n4:thm:lower} 合并，完成主定理~\ref{n4:thm:main} 的证明。
\qed

\subsection{计算证书与复现边界}\label{n4:sec:certificate}

为使计算机辅助部分可审计，证书采用两套彼此独立的实现。
\begingroup
\small

\begin{center}
\begin{tabularx}{0.94\textwidth}{@{}>{\hsize=0.78\hsize}L>{\hsize=1.22\hsize}L@{}}
\toprule
文件 & 作用\\
\midrule
\path{perm4_quadratic_extension_chart_exact_verify.py}
& 在 \(\Q\) 上重建式~\eqref{n4:eq:659}，计算 Schur 补、
\(\det S_0=-32768\)，并核验十五个归一化矩阵同时严格三角化。\\
\path{perm4_quadratic_extension_full_minor_replay.py}
& 在两个独立素数及多个参数点直接重放完整 \(659\times659\) 子式，
只作交叉核验，不承担从有限域到特征零的逻辑跳跃。\\
\path{perm4_rank8_independent_audit.py}
& 不导入项目模块，从显式基重新构造 \(560\)、\(92\)、三次延拓和
\(659\times659\) 图表子式，防止共享辅助函数或索引错误。\\
\path{perm4_rank8_verify_all.py}
& 一条命令运行上界、基础秩、精确图表证书和完整子式重放。\\
\bottomrule
\end{tabularx}
\end{center}

独立审计的关键输出为
\begin{verbatim}
independent_perm3_prolongations=16
independent_chow_koszul_rank=92
independent_chow_triple_prolongations=4
direct_E_image_rank=560
independent_chart_constant=-32768
independent_chart_triangular=True
INDEPENDENT_RANK8_AUDIT_PASS
\end{verbatim}

逻辑依赖边界如下：
\begin{enumerate}[label=(\arabic*)]
  \item 基础秩 \(560\) 与 \(92\) 的下界由明确整数子式模
  \(1000003\) 非零证明；其上界分别由 \(16\) 与 \(4\) 个显式
  三次延拓给出。
  \item 全局图表命题的核心恒等式
  \(\det\mathcal M(v)=-32768v_{00}^{99}\) 在 \(\Q\) 上精确验证。
  \item 从一个图表推广到所有 \(v\neq0\) 只使用
  \(S_4\times S_4\) 的行列置换对称性。
  \item 其余论证均为本文展开的线性代数。因此结论是严格的
  计算机辅助定理，而不是数值猜想。
\end{enumerate}

\paragraph{推荐复现命令。}
在项目目录中运行
\begin{verbatim}
python perm4_rank8_independent_audit.py
python perm4_rank8_verify_all.py
\end{verbatim}
两条命令都必须正常退出，并分别出现
\texttt{INDEPENDENT\_RANK8\_AUDIT\_PASS} 与总体验收通过标志。
\endgroup
